% §5.6 Bibliographic Notes
%
% Cross-reference resolution status:
%
% ALREADY RESOLVED (within chapter 5, drafted in Tasks 2-6 of Plan 5):
%   \Cref{sec:from-atom-to-set}                 -- chapter 5 §5.1
%   \Cref{sec:bloom-from-scratch}               -- chapter 5 §5.2
%   \Cref{sec:predicate-algebra}                -- chapter 5 §5.3
%   \Cref{sec:one-sided-two-sided}              -- chapter 5 §5.4
%   \Cref{def:one-sided}                        -- chapter 5 §5.4
%
% ALREADY RESOLVED (Plan 0 chapter stubs):
%   \Cref{ch:set-composition}                   -- chapter 6 stub label
%   \Cref{ch:optimal-bernoulli-sets}            -- chapter 14 stub label
%   \Cref{ch:perfect-hash-filters}              -- chapter 15 stub label
%
% Citations resolved by new references.bib entries added in this task:
%   \cite{broder-mitzenmacher-2004}             -- Broder and Mitzenmacher 2004
%   \cite{carter-wegman-1979}                   -- Carter and Wegman 1979
%   \cite{towell-bernoulli-hash-function}       -- Towell BHF manuscript
%
% Citations resolved by prior references.bib entries:
%   \cite{bloomfilter}                          -- already in references.bib
%   \cite{mitzenmacher-upfal}                   -- already in references.bib (Plan 3)
%   \cite{towell-bernoulli-sets}                -- already in references.bib (Plan 0)
%   \cite{towell-bernoulli-data-type}           -- already in references.bib (Plan 1)

\section*{Bibliographic Notes}
\addcontentsline{toc}{section}{Bibliographic Notes}
\label{sec:bib-notes-ch5}

\paragraph{Bloom's original paper.}
The data structure at the center of \Cref{sec:bloom-from-scratch} is due
to \textcite{bloomfilter}, in a four-page \emph{Communications of the
ACM} note from 1970. Bloom's framing was operational rather than
probabilistic: he was concerned with hyphenation dictionaries for a
typesetting application, and the question was how much memory could be
saved by accepting a small rate of mishyphenated words in exchange for
not having to store the dictionary itself. The analysis in that paper
already contains the optimal-$k$ derivation that
\Cref{thm:bloom-optimal-k} reproduces in modern notation, including
the $k = (m/n)\ln 2$ result and the resulting per-bit
false-positive rate $(1/2)^k$ at the optimum. What this textbook adds
is not the construction or its analysis but the framing: the Bloom
filter as a particular construction of a function
$\Universe \to \Bernoulli{\BoolType}$, rather than as a clever bit-array
trick for set membership. The two readings agree on every numerical
prediction.

\paragraph{Survey and textbook context.}
For an entry into the broader literature on approximate-set data
structures, two references are standard. \textcite{mitzenmacher-upfal}
devotes Chapter~5 to hashing, Bloom filters, and the
balls-and-bins analysis that underlies both; the proof of the
optimal-$k$ result in that text follows the same calculus argument used
in \Cref{sec:bloom-from-scratch}, though without the
$\Bernoulli{\BoolType}$ framing. \textcite{broder-mitzenmacher-2004}
collects applications of Bloom filters across networking, distributed
systems, and database query processing, organized by the kind of
approximate-set operation each application requires. The survey is
the natural follow-up reading for anyone who wants to see how widely
the Bloom filter's $\fnrate = 0$ property gets used in practice; it
also documents the variants (counting filters, scalable filters,
compressed filters) that the present chapter mentions only in passing.

\paragraph{The set-as-predicate viewpoint.}
The framing of an approximate set as a function
$\Universe \to \Bernoulli{\BoolType}$, rather than as a data structure
that happens to have an error rate, is the contribution of
\textcite{towell-bernoulli-sets}. Earlier literature, including
\textcite{bloomfilter} and the broad sweep of
\textcite{broder-mitzenmacher-2004}, treats approximate sets as
\emph{data structures with error rates}: an array of bits plus a
membership test, characterized by a single false-positive probability.
The Bernoulli-sets reading inverts that emphasis. The object of study
is the function itself, the per-query Boolean output is a
$\Bernoulli{\BoolType}$ instance in the sense of
\textcite{towell-bernoulli-data-type}, and the data structure is one
of many possible realizations. \Cref{sec:from-atom-to-set,sec:predicate-algebra}
develop this framing as a continuation of Part~I's atom-and-axiom
construction; the algebra of \Cref{sec:predicate-algebra} is the
direct lift of the order-2 $\Bernoulli{\BoolType}$ algebra of
chapter~4 to the per-element setting.

\paragraph{One-sided and perfect-filter lineage.}
The one-sided / two-sided dichotomy of \Cref{sec:one-sided-two-sided}
has two ancestors in the literature. The FNR-zero shape is the Bloom
filter's structural property, named as such in
\textcite{bloomfilter} but not isolated as a category until the
broader approximate-set literature accumulated multiple constructions
sharing it. The FPR-zero shape---a \emph{perfect filter} in the sense
of \Cref{def:one-sided}, on a designated set---has its constructive
roots in the perfect-hashing literature opened by
\textcite{carter-wegman-1979}. Carter and Wegman's universal hash
families and the perfect hashing constructions that followed are
what makes a fingerprint-per-element scheme realizable in the first
place: without a perfect hash on $A$, the fingerprint table cannot
guarantee $\fprate = 0$ on $x \in A$. Chapter~15's perfect hash
filter (\Cref{ch:perfect-hash-filters}) is the modern descendant
of that line. The chapter's exposition treats the Bernoulli-set
characterization first, then the construction; the historical order is
the opposite.

\paragraph{Forward bridge: chapters 14 and 15.}
Both chapter~14 (\Cref{ch:optimal-bernoulli-sets}) and chapter~15
(\Cref{ch:perfect-hash-filters}) build on a single foundational
manuscript: \textcite{towell-bernoulli-hash-function}, which derives
the space-optimal Bernoulli hash filter (BHF) for both the unrestricted
and FPR-zero-on-$A$ regimes and proves the entropy lower bound that
makes ``optimal'' precise. The construction in that paper realizes
the FPR-zero shape of \Cref{def:one-sided} explicitly and shows that
the Bloom filter is, on the relevant axis, off-optimal by a constant
factor that is well understood. Readers who want the punch line of
the chapters that follow---why \emph{optimal} can be said precisely,
and what the cost of optimality is in construction time and update
support---will find it in that manuscript.
